%% kapa: exam=1 round=33 year=2022 period=1 form=A subject=real_estate subject_ko=부동산학원론 seq=37 number=117
%% source: Q-Net 감정평가사 기출문제 (cst00309) artlSeq=5209997 file=raw/1008/5209997/2022년도 제33회 감정평가사 1차시험 1교시(민법,경제학,부동산학).pdf
%% conversion: kordoc 4.24.51 -> tools/gen1.py (auto); review=proofread-sonnet; flags=none
\begin{question}{117}
수익환원법에 관한 설명으로 옳지 않은 것은?
\begin{choices}
  \item 운영경비에 감가상각비를 포함시킨 경우 상각전환원율을 적용한다.
  \item 직접환원법에서 사용할 환원율은 시장추출법으로 구하는 것을 원칙으로 한다.
  \item 재매도가치를 내부추계로 구할 때 보유기간 경과 후 초년도 순수익을 반영한다.
  \item 할인 또는 환원할 순수익을 구할 때 자본적지출은 비용으로 고려하지 않는다.
  \item 요소구성법으로 환원율을 결정할 때 위험요소를 적극적으로 반영하면 환원율은 커진다.
\end{choices}
\end{question}
\answer{1}
